\chapter{Introduction}
Modern data applications rarely store all useful information in one representation.
A movie catalogue contains exact attributes such as year, runtime, director, and
genre, while its reviews contain opinions that cannot be recovered reliably with
SQL string matching. The motivating query therefore combines both: ``Which movies
released in 2001 have reviews that explicitly recommend the movie?''

The project objective was to produce machine-readable answer rows from such
questions with three simultaneous goals: high precision and recall, low wall-clock
latency, and few expensive model calls. The initial systems exposed a trade-off.
SUQL produced relatively good answers by executing structured predicates before
row-wise semantic calls, but was slow. A semantic block join placed many rows in a
prompt and supported separate sources, but its answer quality and unnecessary work
on structured mismatches were problematic.

The work progressed from reproducing these baselines to adding two-level model
cascades. SUQL V1 retains row-wise execution and routes uncertain log-odds scores to
a strong model. Trummer V1 composes structured pruning, an exact key join, cheap
batches, and coalesced strong-model batches. I also developed a canonical benchmark,
ground-truth annotations, common runners, model-call instrumentation, and plots.

The main research question is: \emph{Can relational pruning and batched heterogeneous
model routing jointly improve the quality--latency--call trade-off of hybrid semantic
queries?} Subsidiary questions concern how confidence should be defined, how a
threshold can be selected without a labelled calibration set, and which failures
must escalate or fail closed.

The report first reviews related execution models, then formalizes the design,
documents the implementation and prompts, evaluates four systems, and closes with
limitations and a roadmap. All numerical claims identify the exact experiment from
which they originate.
